\begin{tikzpicture}[
  x=1mm,y=1mm,figure picture,
  cell/.style={rectangle,rounded corners=.6mm,draw=#1,fill=white,
    line width=.6pt,minimum width=7.4mm,minimum height=5.4mm,inner sep=0pt,
    font=\figurefont\scriptsize},
  row label/.style={font=\figurefont\scriptsize,text=FigureInk,anchor=east},
  head label/.style={font=\figurefont\scriptsize,text=FigureInk,anchor=west},
  up/.style={draw=FigureMuted,line width=.5pt,-{Stealth[length=1.3mm,width=.9mm]}},
  flow/.style={figure flow}
]
  \def\dx{8.6}
  % 输入词元（第一段为问题，第二段为材料）
  \foreach \tok [count=\i from 0] in {{[CLS]},改,到,哪,天,{[SEP]},会,议,改,至,周,五,{[SEP]}}{
    \node[font=\figurefont\scriptsize] (t\i) at (\i*\dx,0) {\tok};
    \node[cell=FigureInput] (e\i) at (\i*\dx,8) {$E_{\text{\tok}}$};
    \node[cell=FigureModel] (s\i) at (\i*\dx,15)
      {\ifnum\i<6 $E_A$\else $E_B$\fi};
    \node[cell=FigureMuted] (p\i) at (\i*\dx,22) {$E_{\i}$};
    \draw[up] (p\i.north) -- (\i*\dx,29);
    \draw[up] (\i*\dx,39) -- (\i*\dx,43.3);
    \node[cell=FigureOutput] (h\i) at (\i*\dx,46)
      {\ifnum\i=0 $\bm h_{\mathrm{C}}$\else $\bm h_{\i}$\fi};
  }
  \node[row label] at (-5,0) {输入词元};
  \node[row label] at (-5,8) {词元嵌入};
  \node[row label] at (-5,15) {$+$ 段嵌入};
  \node[row label] at (-5,22) {$+$ 位置嵌入};
  \node[row label] at (-5,46) {末层表示};

  % 编码器
  \draw[draw=FigureStroke,fill=FigurePanel,rounded corners=1.5mm,line width=.75pt] (-4.5,29) rectangle (12*\dx+4.5,39);
  \node[font=\figurefont\footnotesize] at (6*\dx,34)
    {$L$层双向Transformer编码器：每个位置的注意力可读取全部13个位置};

  % 分类头
  \node[figure node,draw=FigureOutput,minimum height=6mm]
    (cls) at (1.5*\dx,60) {分类头\\$\operatorname{softmax}(\bm W\bm h_{\mathrm{C}})$};
  \draw[flow] (h0.north) -- (h0.north|-cls.south);

  % 序列标注头：材料各字的标签
  \foreach \i/\tag in {6/O,7/O,8/O,9/O,10/B,11/I}{
    \node[font=\figurefont\scriptsize,text=FigureInk] (g\i) at (\i*\dx,54) {\tag};
  }
  % 跨度头：起点与终点
  \draw[draw=BookTeal,line width=1pt] (10*\dx-3,58) -- (10*\dx-3,61) -- (11*\dx+3,61) -- (11*\dx+3,58);
  \node[font=\figurefont\scriptsize,text=BookTeal] at (10*\dx,64) {起点};
  \node[font=\figurefont\scriptsize,text=BookTeal] at (11*\dx,64) {终点};
  \node[head label] at (12*\dx+5,54) {序列标注头};
  \node[head label] at (12*\dx+5,61) {跨度头};
\end{tikzpicture}
